% 全书图形的基础语法；局部几何与数据编码可覆盖对应参数。
\tikzset{
  figure picture/.style={font=\figurefont\footnotesize,text=FigureInk,
    line width=.6pt},
  figure node/.style={rectangle,rounded corners=1.5mm,
    draw=FigureStroke,fill=FigureModelFill,line width=.75pt,
    minimum height=10mm,inner xsep=3pt,inner ysep=4pt,
    font=\figurefont\footnotesize,align=center,text=FigureInk},
  figure flow/.style={-{Stealth[length=1.8mm,width=1.2mm]},
    draw=FigureStroke,line width=.85pt},
  figure link/.style={draw=FigureMuted,line width=.6pt},
  figure panel/.style={draw=FigureStroke,fill=white,line width=.75pt,
    rounded corners=2.5mm},
  figure subgroup/.style={draw=FigureMuted,line width=.6pt,
    rounded corners=2mm,dash pattern=on 2pt off 2pt},
  figure heading/.style={font=\figurefont\bfseries\footnotesize,text=FigureInk},
  figure note/.style={font=\figurefont\footnotesize,text=FigureMuted,align=center},
  figure edge label/.style={font=\figurefont\footnotesize,text=FigureInk,
    fill=white,inner sep=1.5pt,align=center},
  book node/.style={figure node,minimum height=10mm,inner sep=6pt,
    font=\figurefont\small},
  book arrow/.style={figure flow}
}
% 左下 x/y、宽、高、标题、标题带颜色。仅用于已有的阶段或模块分区。
% 使用 x=y=1mm 的坐标；在内容之前绘制，以免遮挡节点与连线。
\newcommand{\figurePanel}[6]{%
  \path[figure panel] (#1,#2) rectangle ({#1+#3},{#2+#4});%
  \begin{scope}
    \clip[rounded corners=2.5mm] (#1,#2) rectangle ({#1+#3},{#2+#4});%
    \fill[#6] (#1,{#2+#4-8}) rectangle ({#1+#3},{#2+#4});%
  \end{scope}
  \draw[draw=FigureStroke,line width=.6pt]
    (#1,{#2+#4-8})--({#1+#3},{#2+#4-8});%
  \node[figure heading] at ({#1+#3/2},{#2+#4-4}) {#5};%
}

% 其他卷保留既有 book node/book arrow；第一卷局部应用科研样式。
\tikzset{book node/.style={draw=none,fill=BookPaper,minimum height=10mm,
  inner sep=6pt,font=\heiti\small,align=center},
  book arrow/.style={-{Stealth[length=2mm]},line width=0.7pt,draw=BookMuted}}
